\section{Fungsi Hilbert dan dimensi gelanggang lokal}
\subsection{Polinomial bernilai bulat}

Sekarang kita perlu membahas sedikit aljabar umum.

Misalkan $P \in \mathbb{Q}[t]$. Kita meninjau kapan $P$ memetakan
bilangan-bilangan bulat $\mathbb{Z}$, atau lebih umum bilangan-bilangan bulat yang cukup besar, ke dalam
$\mathbb{Z}$. Tentu saja, setiap polinomial dalam $\mathbb{Z}[t]$ mempunyai sifat ini, tetapi
ada pula yang lain; misalnya, tinjau $\frac{1}{2}(t^2 -t)$.

\begin{proposition}\label{integervalued} 
Misalkan $P \in \mathbb{Q}[t]$. Maka $P(m)$ merupakan bilangan bulat untuk setiap bilangan bulat $m \gg 0$ jika dan hanya jika
$P$ dapat dituliskan dalam bentuk
\[ P(t) = \sum_n c_n \binom{t}{n}, \quad c_n \in \mathbb{Z}.  \]
Khususnya, $P(\mathbb{Z}) \subset \mathbb{Z}$.
\end{proposition} 
Jadi $P$ merupakan kombinasi linear-$\mathbb{Z}$ dari koefisien-koefisien binomial. 
\begin{proof} 
Perhatikan bahwa himpunan $\left\{\binom{t}{n}\right\}_{n \in \mathbb{Z}_{\geq 0}}$ membentuk suatu basis bagi ruang
polinomial $\mathbb{Q}[t]$. Karena itu, jelas bahwa $P(t)$ dapat dituliskan sebagai
kombinasi rasional $\sum c_n \binom{t}{n}$ dengan $c_n \in \mathbb{Q}$.
Kita perlu membuktikan bahwa sebenarnya $c_n \in \mathbb{Z}$.

Tinjau operator $\Delta$ yang didefinisikan pada fungsi-fungsi $\mathbb{Z} \to
\mathbb{C}$ sebagai berikut:
\[( \Delta f)(m) = f(m) - f(m-1).  \]
Jelas bahwa jika $f$ bernilai bulat untuk $m \gg 0$, maka demikian pula
$\Delta f$. Mudah pula diperiksa bahwa $\Delta \binom{t}{n} =
\binom{t}{n-1}$.

Dengan meninjau
fungsi $\Delta P = \sum c_n \binom{t}{n-1}$ (yang bernilai dalam $\mathbb{Z}$), mudah dilihat bahwa $c_n \in \mathbb{Z}$ melalui induksi
pada derajat. 
Mudah pula dilihat secara langsung bahwa koefisien binomial bernilai dalam
$\mathbb{Z}$ pada \emph{semua} argumen.
\end{proof} 



\subsection{Definisi dan contoh}
Misalkan $R$ suatu gelanggang. 
\begin{question} 
Apakah definisi yang baik untuk $\dim(R)$? Bahkan, secara lebih umum, 
apakah dimensi $R$ pada suatu ``titik'' tertentu (yakni ideal prima)?
\end{question} 

Secara geometris, bayangkan $\spec R$ untuk sebarang gelanggang; pilih suatu titik yang bersesuaian dengan ideal
maksimal $\mathfrak{m} \subset R$. Kita ingin mendefinisikan \textbf{dimensi $R$}
pada $\mathfrak{m}$. Gagasan ini dapat dibayangkan menyerupai ``dimensi atas
bilangan kompleks'' bagi varietas aljabar yang didefinisikan atas $\mathbb{C}$. Namun definisi itu
harus murni aljabar. 
Apa yang dapat kita lakukan?


Berikut suatu gagasan. Agar ruang topologis $X$ berdimensi $n$ pada $x \in
X$, seharusnya terdapat $n$ koordinat pada titik $x$. Dengan kata lain,
titik $x$ seharusnya ditentukan secara tunggal
oleh lokus nol dari $n$ koordinat pada ruang tersebut.
Termotivasi oleh hal ini, kita dapat mencoba mendefinisikan $\dim_{\mathfrak{m}} R$
sebagai banyaknya pembangkit $\mathfrak{m}$.
Namun definisi ini buruk, sebab $\mathfrak{m}$ mungkin tidak mempunyai jumlah
pembangkit yang sama dengan $\mathfrak{m}R_{\mathfrak{m}}$. Dengan kata lain, ini bukan definisi yang benar-benar
\emph{lokal}.
\begin{example} 
Misalkan $R$ suatu daerah integral Noether yang tertutup secara integral tetapi bukan UFD. Misalkan $\mathfrak{p}
\subset R$ suatu ideal prima yang minimal di atas suatu ideal utama, tetapi
bukan ideal utama. Maka $\mathfrak{p}R_{\mathfrak{p}}$ dibangkitkan oleh
satu unsur, sebagaimana kelak akan kita lihat, sedangkan $\mathfrak{p}$ tidak.
\end{example} 

Kita menginginkan definisi dimensi yang bersifat
lokal.
Hal ini membawa kita pada definisi berikut.
\begin{definition} 
Jika $R$ suatu gelanggang \emph{lokal} (Noether) dengan ideal maksimal $\mathfrak{m}$,
maka \textbf{dimensi pembenaman} dari $R$, yang dilambangkan $\emdim R$, adalah
jumlah minimal pembangkit $\mathfrak{m}$. Jika $R$ suatu gelanggang
Noether dan $\mathfrak{p} \subset R$ suatu ideal prima, maka \textbf{dimensi pembenaman
pada $\mathfrak{p}$} adalah dimensi pembenaman gelanggang lokal $R_{\mathfrak{p}}$.
\end{definition} 

Dalam definisi di atas, jelas bahwa cukup mempelajari apa yang terjadi untuk
gelanggang lokal; untuk sementara kita memberlakukan pembatasan tersebut. Menurut Lema Nakayama,
dimensi pembenaman adalah jumlah minimal pembangkit
$\mathfrak{m}/\mathfrak{m}^2$, atau dimensi atas $R/\mathfrak{m}$ dari ruang vektor
tersebut:
\[ \emdim R = \dim_{R/\mathfrak{m}} \mathfrak{m}/ \mathfrak{m}^2.  \]


Namun secara umum, dimensi pembenaman tidak sama dengan
dimensi ``geometris'' intuitif suatu varietas
aljabar.

\begin{example} 
Misalkan $R = \mathbb{C}[t^2, t^3] \subset \mathbb{C}[t]$, yaitu gelanggang koordinat
kurva kubik $y^2 =x^3$, sebab $R \simeq \mathbb{C}[x,y]/(x^2 - y^3)$
melalui $x = t^3, y = t^2$. Mari kita lokalkan pada ideal prima $\mathfrak{p} = (t^2,
t^3)$; kita memperoleh $R_{\mathfrak{p}}$. 

Spektrum $\spec R$ singular di titik asal. Sebagai akibatnya, $\mathfrak{p}
R_{\mathfrak{p}} \subset R_{\mathfrak{p}}$ memerlukan dua pembangkit, tetapi
varietas yang bersesuaian berdimensi satu.
\end{example} 

Jadi dimensi pembenaman adalah dimensi terkecil dari ruang mulus tempat kita dapat membenamkan
$R$.
Namun bagi varietas singular, ini bukan dimensi yang kita inginkan. 

Alih-alih hanya meninjau $\mathfrak{m}/\mathfrak{m}^2$, mari kita
tinjau \emph{barisan} ruang vektor berdimensi hingga
\[ \mathfrak{m}^k/\mathfrak{m}^{k+1}.  \]
Menghitung dimensi-dimensi ini sebagai fungsi $k$ menghasilkan suatu invarian yang mendeskripsikan geometri
lokal dari $\spec R$.

Kelak akan kita buktikan:
\begin{theorem} \label{hilbfnispolynomial}
Misalkan $(R, \mathfrak{m})$ suatu gelanggang lokal Noether. Maka terdapat suatu
polinomial
$f \in \mathbb{Q}[t]$ sedemikian sehingga
\[ f(n) =  \ell(R/\mathfrak{m}^n) = \sum_{i=0}^{n-1} \dim
\mathfrak{m}^i/\mathfrak{m}^{i+1} \quad \forall n \gg 0.  \]

Selain itu, $\deg f \leq \dim \mathfrak{m}/\mathfrak{m}^2$.
\end{theorem}


Perhatikan bahwa polinomial ini terdefinisi dengan baik, sebab dua polinomial yang sama untuk $n$ besar
pastilah identik. Perhatikan pula bahwa $R/\mathfrak{m}^n$ bersifat Artin sehingga mempunyai panjang hingga,
dan kita telah menggunakan fakta bahwa panjang bersifat aditif pada barisan eksak
pendek. Kita ingin menuliskan $\dim R/\mathfrak{m}^n$, tetapi secara umum
hal itu tidak dapat dilakukan; sebagai gantinya kita menggunakan panjang. 

Berdasarkan hal ini, kita mendefinisikan:
\begin{definition} 
\textbf{Dimensi} gelanggang lokal $R$ adalah derajat polinomial
$f$ di atas. Untuk gelanggang Noether sebarang $R$, kita definisikan $\dim R =
\sup_{\mathfrak{p} \in \spec R} \dim (R_{\mathfrak{p}})$.
\end{definition} 


Sekarang mari kita lakukan beberapa perhitungan contoh. 

\begin{example}[Garis afin] \label{easydimcomputation}
\label{dimaffineline}
Tinjau gelanggang lokal $(R, \mathfrak{m}) = \mathbb{C}[t]_{(t)}$. Maka $\mathfrak{m} = (t)$ dan
$\mathfrak{m}^k/\mathfrak{m}^{k+1}$ berdimensi satu, dibangkitkan oleh $t^k$. Khususnya,
gelanggang tersebut berdimensi satu. 
\end{example} 

\begin{example}[Kurva singular] Tinjau $R = \mathbb{C}[t^2, t^3]_{(t^2, t^3)}$, gelanggang lokal dari $y^2 = x^3$
di nol. Maka $\mathfrak{m}^n$ dibangkitkan oleh $t^{2n}, t^{2n+1}, \dots$.
$\mathfrak{m}^{n+1}$ dibangkitkan oleh $t^{2n+2}, t^{2n+3}, \dots$. Jadi semua
hasil bagi tersebut berdimensi dua. Dimensinya sedikit
lebih besar daripada dalam \rref{easydimcomputation}, tetapi tidak bertumbuh. Gelanggang tersebut tetap berdimensi satu. 
\end{example} 

\begin{example}[Bidang afin] \label{dimaffineplane}
Tinjau $R = \mathbb{C}[x,y]_{(x,y)}$. Maka $\mathfrak{m}^k$ dibangkitkan oleh
polinomial-polinomial dalam $x,y$ yang homogen berderajat $k$. Jadi $\mathfrak{m}^k/\mathfrak{m}^{k+1}$
mempunyai dimensi yang \emph{bertumbuh} secara linear dalam $k$. Ini benar-benar merupakan contoh
berdimensi dua. 
\end{example} 

Perbedaan antara pertumbuhan konstan, linear, dan kuadratik dalam
$R/\mathfrak{m}^n$ ketika $n \to \infty$, bukan ukuran suku-suku awalnya,
itulah yang kita kehendaki dari definisi dimensi.

Sekarang mari kita perumum \rref{dimaffineline} dan \rref{dimaffineplane}
di atas ke ruang afin berdimensi sebarang. 
\begin{example}[Ruang afin] 
Tinjau $R = \mathbb{C}[x_1, \dots, x_n]_{(x_1, \dots, x_n)}$. 
Secara geometris, gelanggang ini merepresentasikan varietas $\mathbb{C}^n = \mathbb{A}^n_{\mathbb{C}}$ di sekitar titik asal, sehingga
secara intuitif seharusnya berdimensi $n$. Mari kita periksa hal ini. 

Kita perlu menghitung polinomial $f$ di atas. Di sini $R/\mathfrak{m}^k$ menyerupai himpunan
polinomial berderajat $<k$ atas $\mathbb{C}$. Dimensinya sebagai ruang vektor
dari ruang ini
diberikan oleh suatu koefisien binomial $\binom{n+k-1}{n}$. Ini merupakan polinomial dalam
$k$ berderajat $n$. Khususnya, $\ell(R/\mathfrak{m}^k)$ bertumbuh seperti $k^n$.
Jadi $R$ berdimensi $n$. 
\end{example} 


Terakhir, kita berikan satu contoh lagi yang menunjukkan bahwa DVR berdimensi satu. Bahkan,
di antara daerah integral lokal Noether yang tertutup secara integral, DVR
\emph{dicirikan} oleh sifat ini (karakterisasi DVR dibuktikan dalam bab mengenai gelanggang Dedekind pada sumber lengkap; pembuktian tersebut berada di luar suplemen terpilih ini).

\begin{example}[Dimensi suatu DVR]
Misalkan $R$ suatu DVR. Maka $\mathfrak{m}^k/\mathfrak{m}^{k+1}$ mempunyai panjang satu untuk
setiap $k$. Jadi $R/\mathfrak{m}^k$ mempunyai panjang $k$. Dengan demikian kita dapat mengambil $f(t) = t$, sehingga
$R$ berdimensi satu.
\end{example} 

\subsection{Fungsi Hilbert merupakan suatu polinomial} 

Meskipun kita telah memberikan definisi dimensi dan menghitung berbagai contoh,
kita belum memeriksa bahwa definisi tersebut terdefinisi dengan baik. 
Kita harus membuktikan \rref{hilbfnispolynomial}.

\begin{proof}[Bukti \rref{hilbfnispolynomial}]
Tetapkan suatu gelanggang lokal Noether $(R, \mathfrak{m})$. Kita akan menunjukkan bahwa
$\ell(R/\mathfrak{m}^n)$ merupakan polinomial untuk $n \gg 0$. Kita juga harus membatasi
derajatnya oleh $\dim_{R/\mathfrak{m}} \mathfrak{m}/\mathfrak{m}^2$, yaitu
dimensi pembenaman. Kita akan melakukannya dengan mereduksi pada suatu fakta umum mengenai
modul bergradasi atas gelanggang polinomial.

Misalkan $S = \bigoplus_n  \mathfrak{m}^n/\mathfrak{m}^{n+1}$. Maka $S$ mempunyai
gradasi alami, dan bahkan secara alami merupakan gelanggang bergradasi melalui
pemetaan perkalian 
\[ \mathfrak{m}^{n_1} \times \mathfrak{m}^{n_2} \to \mathfrak{m}^{n_1 + n_2}.  \]
Sebenarnya, $S$ adalah \emph{gelanggang bergradasi terkait} dari filtrasi $\mathfrak{m}$-adik.
Perhatikan bahwa $S_0 = R/\mathfrak{m}$ merupakan suatu medan, yang akan kita lambangkan dengan $k$.
Jadi $S$ merupakan aljabar-$k$ bergradasi.

\begin{lemma} 
$S$ merupakan aljabar-$k$ terbangkitkan hingga. Bahkan, $S$ dapat dibangkitkan oleh paling banyak
$\emdim(R)$ unsur.
\end{lemma} 
\begin{proof} 
Misalkan $x_1, \dots, x_r$ pembangkit-pembangkit $\mathfrak{m} $ dengan $r = \emdim(R)$.
Unsur-unsur ini (atau lebih tepatnya, citra-citranya) membentuk suatu basis-$k$
bagi $\mathfrak{m}/\mathfrak{m}^2$.
Maka citra-citra mereka di $\mathfrak{m}/\mathfrak{m}^2 \subset S$ membangkitkan $S$.
Hal ini mengikuti karena $S_1$ membangkitkan $S$ sebagai aljabar-$S_0$: hasil kali
unsur-unsur dalam $\mathfrak{m}$ membangkitkan pangkat-pangkat $\mathfrak{m}$ yang lebih tinggi.
\end{proof} 

Jadi $S$ merupakan hasil bagi bergradasi dari gelanggang polinomial $k[t_1, \dots, t_r]$, dengan
$t_i $ dipetakan ke $x_i$. Khususnya, $S$ merupakan modul bergradasi terbangkitkan hingga atas $k[t_1,
\dots, t_r]$.
Perhatikan pula bahwa $\ell(R/\mathfrak{m}^n)  = \dim_{k}(S_0) + \dots +
\dim_{k}(S_{n-1})$ untuk setiap $n$, berkat filtrasi tersebut. Inilah
invarian yang kita minati.

Sekarang cukup dibuktikan proposisi yang lebih umum berikut.
\begin{proposition} \label{hilbfngeneral}
Misalkan $M$ sebarang modul bergradasi terbangkitkan hingga atas gelanggang polinomial
$k[x_1, \dots, x_r]$. Maka
terdapat polinomial $f_M^+ \in \mathbb{Q}[t]$ berderajat $ \leq r$, sedemikian sehingga
\[ f_M^+(t) = \sum_{s \leq t} \dim M_s \quad t \gg 0.   \]
\end{proposition} 
Menerapkannya pada $M = S$ memberikan hasil yang diinginkan. Kita dapat melupakan
semua hal lainnya dan hanya meninjau masalah ini untuk gelanggang polinomial bergradasi.

Fungsi ini disebut \textbf{fungsi Hilbert}.

\begin{proof}[Bukti \rref{hilbfngeneral}]
Perhatikan bahwa jika kita mempunyai barisan eksak modul bergradasi atas gelanggang
polinomial,
\[ 0 \to M' \to M \to M'' \to 0,   \]
dan polinomial $f_{M'}, f_{M''}$ seperti dalam proposisi, maka $f_M$ ada dan 
\[ f_M = f_{M'} + f_{M''}.  \] Hal ini langsung dari definisi.
Selanjutnya, kita amati bahwa jika $M$ merupakan modul bergradasi terbangkitkan hingga atas dua
gelanggang polinomial berbeda, tetapi dengan gradasi yang sama, maka keberadaan (dan
nilai) $f_M$ tidak bergantung pada gelanggang polinomial mana yang ditinjau.
Terakhir, kita amati bahwa cukup dibuktikan bahwa $f_M(t) = \dim M_t$ merupakan
polinomial dalam $t$ untuk $t \gg 0$.

Kita akan
menggunakan ketiga pengamatan ini dan melakukan induksi pada $n$.

Jika $n  = 0$, maka $M$ merupakan ruang vektor bergradasi berdimensi hingga atas
$k$, dan gradasinya harus terkonsentrasi pada berhingga banyaknya derajat. Jadi
hasilnya jelas, sebab $f_{M}(t)$ sama dengan $\dim M$ (yang merupakan
dimensi yang sesuai untuk $t \gg 0$). 

Andaikan $n > 0$. Tinjau
filtrasi $M$
\[ 0 \subset \ker( x_1: M \to M) \subset \ker (x_1^2: M \to M) \subset \dots
\subset M.  \]
Karena sifat Noether, filtrasi ini harus stabil pada suatu $M' \subset M$. Setiap
hasil bagi $\ker( x_1^i)/\ker (x_1^{i+1})$ merupakan modul terbangkitkan hingga atas
$k[x_1, \dots,
x_n]/(x_1)$, yaitu gelanggang polinomial yang lebih kecil. Jadi setiap hasil bagi tersebut
$\ker (x_1^{i+1})/\ker (x_1^{i})$ mempunyai fungsi Hilbert berderajat $\leq n-1$ menurut
hipotesis induksi. 

Dengan menelusuri filtrasi tersebut, kita melihat bahwa $M'$ mempunyai fungsi Hilbert yang merupakan jumlah fungsi Hilbert
dari hasil-hasil bagi $\ker(x_1^{i+1})/\ker(x_1^{i})$. Khususnya, $f_{M'}$ ada. Jika kita menunjukkan bahwa $f_{M/M'}$
ada, maka $f_M$ pasti ada. Jadi cukup ditunjukkan bahwa
fungsi Hilbert $f_M$ ada ketika $x_1$ bukan pembagi nol pada $M$. 

Jadi kita telah mereduksi ke kasus ketika $M \stackrel{x_1}{\to} M$
injektif.
Sekarang $M$ mempunyai filtrasi
\[ M \supset x_1 M \supset x_1^2 M \supset \dots  \]
yang merupakan filtrasi menyeluruh dari $M$, dalam arti bahwa tidak ada unsur yang dapat terus-menerus habis dibagi oleh
pangkat-pangkat $x_1$, sebab jika demikian derajatnya tidak akan berhingga. Jadi
$\bigcap x_1^m M = 0$. 

Misalkan $N = M/x_1 M $, yang isomorfik dengan $ x_1^m M/x_1^{m+1} M$ karena $M \stackrel{x_1}{\to} M$
injektif. Di sini $N$ merupakan modul bergradasi terbangkitkan hingga atas $k[x_2, \dots, x_n]$, dan menurut
hipotesis induksi pada $n$, terdapat polinomial $f_N^+$ berderajat $\leq n-1$ sedemikian sehingga
\[ f_N^+(t) = \sum_{t' \leq t} \dim N_{t'}, \quad t \gg 0.  \]

Tetapkan $t \gg 0$ dan tinjau ruang vektor-$k$ $M_t$, yang mempunyai filtrasi hingga 
\[ M_t \supset (x_1 M)_t \supset (x_1^2 M)_t \supset \dots  \]
dengan hasil-hasil bagi berturut-turut berupa komponen bergradasi dari $N \simeq
M/x_1 M \simeq x_1 M/x_1^2 M \simeq \dots$ pada dimensi $t, t-1, \dots$. Kita
memperoleh, misalnya,
\[ (x_1^2 M)_t/(x_1^3 M)_t \simeq N_{t-2},  \]
Dengan menjumlahkannya, kita memperoleh
\[ \dim M_t = \dim N_t + \dim N_{t-1} + \dots . \]
Jumlah di atas sebenarnya berhingga. Memang, dari keterbangkitan hingga terdapat $K
\gg 0 $ sedemikian sehingga $\dim N_q  = 0$ untuk $q< -K$. Dari sini diperoleh 
\[ \dim M_t = \sum_{t' = -K}^{t} \dim N_{t'},  \]
yang mengimplikasikan bahwa $\dim M_t$ merupakan polinomial untuk $t \gg 0$. Ini menyelesaikan
pembuktian.
\end{proof} 
\end{proof}


Misalkan $(R, \mathfrak{m})$ suatu gelanggang lokal Noether dan $M$ suatu modul-$R$
terbangkitkan hingga. 
\begin{proposition} \label{hilbertlocalring} 
$\ell(M/\mathfrak{m}^m M)$ merupakan polinomial untuk $m \gg 0$. 
\end{proposition} 
\begin{proof} 
Pernyataan ini mengikuti dari \rref{hilbfngeneral}; bahkan, pada dasarnya kita telah
melihat argumennya di atas. Tinjau
modul bergradasi terkait
\[ N = \bigoplus \mathfrak{m}^k M/\mathfrak{m}^{k+1}M , \]
yang terbangkitkan hingga atas gelanggang bergradasi terkait
\[ \bigoplus \mathfrak{m}^k/\mathfrak{m}^{k+1}.  \]
Akibatnya, dimensi komponen-komponen bergradasi $N$ bertumbuh secara polinomial
untuk derajat yang besar. Ini mengimplikasikan hasil tersebut. 

\end{proof} 
\begin{definition} 
Kita definisikan \textbf{fungsi Hilbert} $H_M(m)$ sebagai polinomial tunggal
sedemikian sehingga 
\[ H_M(m) = \ell(M/ \mathfrak{m}^m M), \quad m \gg 0.  \]
\end{definition} 

Kebetulan, jelas bahwa $H_M$ bernilai bulat; jadi menurut
\rref{integervalued}, $H_M$ merupakan kombinasi linear-$\mathbb{Z}$
dari koefisien-koefisien binomial.

\subsection{Dimensi suatu modul}
Misalkan $R $ suatu gelanggang lokal Noether dengan ideal maksimal $\mathfrak{m}$. Kita
telah melihat (\rref{hilbertlocalring}) bahwa terdapat polinomial $H(t)$ dengan
\[ H(t) = \ell(R/\mathfrak{m}^t), \quad t \gg 0.  \]
Sebelumnya, kita mendefinisikan \textbf{dimensi} $R$ sebagai derajat $f_M^+$. 
Karena derajat fungsi Hilbert paling besar sama dengan jumlah pembangkit
gelanggang polinomial, kita telah melihat bahwa 
\[ \dim R \leq \emdim R.  \]

Dengan perangkat fungsi Hilbert, kita dapat memperluas
definisi ini ke modul.
\begin{definition} 
Jika $R$ lokal Noether dan $N$ suatu modul-$R$ hingga, maka $N$ mempunyai
polinomial Hilbert $H_N(t)$ yang, ketika dievaluasi pada $t \gg 0$, memberikan
panjang $\ell(N/\mathfrak{m}^t N)$.
Kita menyebut \textbf{dimensi
$N$} sebagai derajat polinomial Hilbert ini.
\end{definition} 

Jelas bahwa dimensi \emph{gelanggang} $R$ sama dengan dimensi
\emph{modul} $R$.

Berikutnya kita menunjukkan bahwa dimensi berperilaku baik terhadap barisan eksak
pendek. Hal ini sebenarnya agak halus sebab, secara umum, penensoran dengan
$R/\mathfrak{m}^t$ tidak eksak; ternyata, menurut Lema Artin--Rees, operasi itu \emph{hampir}
eksak. 
Di sisi lain, fakta yang bersesuaian untuk modul atas suatu \emph{gelanggang
polinomial} sangat mudah, sebab 
definisinya tidak melibatkan penensoran.

\begin{proposition} 
Andaikan kita mempunyai barisan eksak
\[ 0 \to M' \to M \to M'' \to 0 \]
modul bergradasi atas gelanggang polinomial $k[x_1, \dots, x_n]$.	Maka
\[ f_M(t) = f_{M'}(t) + f_{M''}(t), \quad  f_M^+(t) = f_{M'}^+(t) +
f_{M''}^+(t). \]
Akibatnya, $\deg f_M = \max \deg f_{M'}, \deg f_{M''}$.
\end{proposition} 
\begin{proof} Bagian pertama langsung karena dimensi bersifat aditif pada ruang
vektor. Bagian kedua mengikuti karena fungsi Hilbert mempunyai koefisien utama
tak negatif.
\end{proof} 

\begin{proposition} \label{dimexactseq}
Tetapkan suatu barisan eksak
\[ 0 \to N' \to N \to N'' \to 0  \]
modul-$R$ hingga. Maka $\dim N = \max (\dim N', \dim N'')$.
\end{proposition} 
\begin{proof} 
Kita mempunyai barisan eksak
\[ 0 \to K \to N/\mathfrak{m}^t N \to N''/\mathfrak{m}^t N'' \to 0  \]
dengan $K$ sebagai kernel. Di sini $K = (N' + \mathfrak{m}^t N)/ \mathfrak{m}^t N
= N'/( N' \cap \mathfrak{m}^t N)$. Ini tidak persis sama dengan $N'/\mathfrak{m}^t N'$,
tetapi cukup dekat. 
Kita mempunyai surjeksi
\[ N'/\mathfrak{m}^t N \twoheadrightarrow N'/(N' \cap \mathfrak{m}^t N) = K. \]
Khususnya, 
\[ \ell(K) \leq \ell(N'/\mathfrak{m}^t N').  \]
Di sisi lain, Lema Artin--Rees memberikan pertidaksamaan dalam
arah sebaliknya. Kita mempunyai inklusi
\[ \mathfrak{m}^t N' \subset N' \cap \mathfrak{m}^t N \subset
\mathfrak{m}^{t-c} N'  \]
untuk suatu $c$. Ini mengimplikasikan bahwa $\ell(K) \geq \ell( N'/\mathfrak{m}^{t-c} N')$. 

Definisikan $M = \bigoplus \mathfrak{m}^t N/\mathfrak{m}^{t+1} N$, dan definisikan $M',
M''$ dengan cara serupa menggunakan $N', N''$. Maka kita telah melihat bahwa 
\[ \boxed{f_M^+(t-c) \leq \ell(K) \leq f_M^+(t).}  \]
Kita juga mengetahui bahwa panjang $K$ ditambah panjang $N''/\mathfrak{m}^t N''$
adalah $f_M^+(t)$, yakni
\[ \ell(K) + f_{M''}^+(t) = f_M^+(t).  \]
Sekarang panjang $K$ merupakan polinomial dalam $t$ yang sangat mirip dengan
$f_{M'}^+$, yaitu koefisien utamanya sama. Jadi kita mempunyai
persamaan hampiran $f_{M'}^+(t) + f_{M''}^+(t) \simeq f_M^+(t)$. Ini mengimplikasikan
hasil tersebut karena derajat $f_M^+$ adalah $\dim N$ (dan demikian pula untuk yang lain). 
\end{proof} 



\begin{proposition} 
$\dim R$ sama dengan $\dim R/\rad R$.
\end{proposition} 
Dengan kata lain, dimensi tidak berubah ketika unsur-unsur nilpoten dihilangkan, sebagaimana
yang kita harapkan, sebab unsur nilpoten tidak memengaruhi $\spec (R)$.
\begin{proof} 
Untuk itu, kita memerlukan sedikit informasi tambahan mengenai fungsi Hilbert.
Karena itu, kita akan menyimpang cukup jauh. 


Terakhir, mari kita kembali ke pernyataan mengenai dimensi dan unsur nilpoten. Misalkan $R$
suatu gelanggang lokal Noether dan $I = \rad (R)$. Maka $I$ suatu modul-$R$ hingga. Khususnya,
$I$ nilpoten, sehingga $I^n  = 0$ untuk $n \gg 0$. Kita akan menunjukkan bahwa
\[ \dim R/I = \dim R/I^2 = \dots \]
yang akan mengimplikasikan hasil tersebut, sebab akhirnya pangkat-pangkat itu menjadi nol. 

Khususnya, untuk setiap $k$ kita harus menunjukkan 
\[ \dim R/I^k  = \dim R/I^{k+1}.  \]
Terdapat barisan eksak
\[ 0 \to I^k/I^{k+1} \to R/I^{k+1} \to R/I^k \to 0.  \]
Dimensi gelanggang-gelanggang ini sama dengan dimensinya sebagai
modul-$R$. Jadi kita dapat menggunakan barisan eksak pendek modul ini. Berdasarkan
hasil sebelumnya, kita direduksi untuk menunjukkan bahwa
\[ \dim I^k/I^{k+1} \leq \dim R/I^k.  \]
Perhatikan bahwa $I$ menganihilasi $I^k/I^{k+1}$. Khususnya, $I^k/I^{k+1}$ merupakan modul
terbangkitkan hingga atas $R/I^k$. Terdapat barisan eksak
\[ \bigoplus_N R/I^k \to I^k/I^{k+1} \to 0  \]
yang mengimplikasikan bahwa $\dim I^k/I^{k+1} \leq \dim \bigoplus_N R/I^k = \dim R/I^k$.
\end{proof} 

\begin{example} 
Misalkan $\mathfrak{p} \subset \mathbb{C}[x_1, \dots, x_n]$ dan misalkan $R =
(\mathbb{C}[x_1,\dots, x_n]/\mathfrak{p})_{\mathfrak{m}}$ untuk suatu ideal maksimal
$\mathfrak{m}$. Berapakah $\dim R$? 
Apakah arti dimensi bagi gelanggang koordinat atas $\mathbb{C}$?

Ingat dari Teorema normalisasi Noether bahwa terdapat suatu gelanggang polinomial
$\mathbb{C}[y_1, \dots, y_m]$ yang termuat dalam $S=\mathbb{C}[x_1,\dots,
x_n]/\mathfrak{p}$, dan $S$ merupakan perluasan integral hingga atas gelanggang polinomial ini. 
Kita klaim bahwa 
\[ \dim R = m.  \]
Hari ini tidak ada cukup waktu untuk membuktikannya. 
\end{example} 

\subsection{Dimensi hanya bergantung pada dukungan}

Misalkan $(R, \mathfrak{m})$ suatu gelanggang lokal Noether. Misalkan $M$ suatu modul-$R$
terbangkitkan hingga.	Kita mendefinisikan \textbf{polinomial Hilbert} dari $M$ sebagai
polinomial yang pada $t \gg 0$ bernilai $\ell(M/\mathfrak{m}^tM)$. Sebelumnya kita membuktikan
bahwa polinomial seperti itu selalu ada, dan menyebut derajatnya sebagai
\textbf{dimensi $M$}. Namun, 
sekarang akan kita lihat bahwa $\dim M$ sebenarnya hanya bergantung pada dukungan\footnote{
Ingat bahwa $\supp M = \left\{\mathfrak{p}: M_{\mathfrak{p}}\neq 0\right\}$.} $\supp M$.
Dalam pengertian ini, dimensi sesungguhnya merupakan pernyataan mengenai \emph{ruang
topologis} $\supp M \subset \spec R$, bukan mengenai $M$ itu sendiri. 


Dengan kata lain, kita akan membuktikan:
\begin{proposition} 
$\dim M$ hanya bergantung pada $\supp M$.
\end{proposition} 

Bahkan, kita akan menunjukkan:

\begin{proposition} 
$\dim M = \max_{\mathfrak{p} \in \supp M} \dim R/\mathfrak{p}$.
\end{proposition} 
\begin{proof} 
Menurut \rref{filtrationlemma} dalam \rref{noetherian}, terdapat suatu filtrasi hingga 
\[ 0 = M_0 \subset M_1 \subset \dots \subset M_m = M,  \]
sedemikian sehingga setiap hasil bagi berturut-turut isomorfik dengan $R/\mathfrak{p}_i
\subset R$
untuk suatu ideal prima $\mathfrak{p}_i$. Diberikan suatu barisan eksak pendek
modul, kita mengetahui bahwa dimensi suku tengah adalah maksimum dimensi kedua
ujungnya (\rref{dimexactseq}). Dengan mengiterasikan ini, kita melihat bahwa
dimensi $M$ adalah maksimum dari
dimensi hasil-hasil bagi berturut-turut $M_i/M_{i-1}$.

Namun semua $\mathfrak{p}_i$ yang muncul
berada dalam $\supp M$, sehingga diperoleh 
\[ \dim M = \max_{\mathfrak{p}_i} R/\mathfrak{p}_i \leq \max_{\mathfrak{p} \in \supp M} \dim R/\mathfrak{p}.  \]
Kita harus menunjukkan pertidaksamaan sebaliknya. Tetapkan sebarang ideal prima $\mathfrak{p} \in \supp
M$. Maka $M_{\mathfrak{p}} \neq 0$, sehingga salah satu $R/\mathfrak{p}_i$, setelah dilokalkan
pada $\mathfrak{p}$, harus tak nol, sebab pelokalan merupakan funktor eksak. Jadi
$\mathfrak{p}$ harus memuat suatu $\mathfrak{p}_i$. Maka $R/\mathfrak{p}$ merupakan
hasil bagi dari $R/\mathfrak{p}_i$. Khususnya,
\[ \dim R/\mathfrak{p} \leq \dim R/\mathfrak{p}_i.  \]
\end{proof} 

Setelah membuktikannya, kita tidak lagi menggunakan notasi $\dim M$, dan selanjutnya menuliskan
$\dim \supp M$.



%% Catatan: materi berikut mengenai gelanggang afin seharusnya ditambahkan dalam 
\begin{comment}
\subsection{Dimensi gelanggang afin} Sebelumnya kita membuat suatu klaim. Jika $R$
suatu daerah integral dan modul hingga atas gelanggang polinomial $k[x_1, \dots, x_n]$,
maka $R_{\mathfrak{m}}$ untuk setiap ideal maksimal $\mathfrak{m} \subset R$ berdimensi
$n$. Hal ini menghubungkan dimensi dengan derajat transendensi. 

Pertama, mari kita bahas perluasan hingga gelanggang. Misalkan $R$ suatu gelanggang komutatif
dan $R \to R'$ suatu morfisme yang menjadikan $R'$ modul-$R$ terbangkitkan hingga (khususnya,
integral atas $R$). Misalkan $\mathfrak{m}' \subset R'$ maksimal. Misalkan
$\mathfrak{m}$ prapetanya di $R$, yang juga maksimal (sebab $R \to R'$
integral). 
Misalkan $M$ suatu modul-$R'$ terbangkitkan hingga, sehingga juga merupakan modul-$R$ terbangkitkan hingga. 

Kita dapat memandang $M_{\mathfrak{m}}$ sebagai modul-$R_{\mathfrak{m}}$ atau
$M_{\mathfrak{m}'}$ sebagai modul-$R'_{\mathfrak{m}'}$. Keduanya
terbangkitkan hingga. 

\begin{proposition} 
$\dim \supp M_{\mathfrak{m}}  \geq \dim \supp M_{\mathfrak{m}'}$.
\end{proposition} 
Di sini $M_{\mathfrak{m}}$ merupakan modul-$R_{\mathfrak{m}}$, sedangkan $M_{\mathfrak{m}'}$
merupakan modul-$R'_{\mathfrak{m}'}$. 

\begin{proof} 
Tinjau $R/\mathfrak{m} \to R'/\mathfrak{m} R' \to R'/\mathfrak{m}'$. Kita
melihat bahwa $R'/\mathfrak{m} R'$ merupakan modul-$R/\mathfrak{m}$ hingga, sehingga menjadi
ruang vektor-$R/\mathfrak{m}$ berdimensi hingga. Khususnya,
$R'/\mathfrak{m} R'$ mempunyai panjang hingga sebagai modul-$R/\mathfrak{m}$, dan
khususnya merupakan gelanggang Artin. Jadi gelanggang ini merupakan hasil kali gelanggang-gelanggang lokal Artin.
Gelanggang-gelanggang Artin ini adalah pelokalan $R'/\mathfrak{m}R'$ pada ideal-ideal
$R'$ yang terletak di atas $\mathfrak{m}$. Salah satu ideal tersebut adalah $\mathfrak{m}'$. 
Jadi khususnya,
\[ R'/\mathfrak{m}R \simeq R'/\mathfrak{m}'\times \mathrm{faktor \ lain}.  \]
Karena nilradikal suatu gelanggang Artin nilpoten, kita melihat bahwa
$\mathfrak{m}'^c R'_{\mathfrak{m}'} \subset \mathfrak{m} R'_{\mathfrak{m}}$ untuk
suatu $c$.

Baiklah, saya tidak dapat mengikuti ini---terlalu lelah. Akan saya lanjutkan suatu hari nanti.
\end{proof} 


\begin{proposition} 
$\dim \supp M_{\mathfrak{m}} = \max_{\mathfrak{m}' \mid \mathfrak{m}} \dim
\supp M_{\mathfrak{m}'}$.
\end{proposition} 

Ini berarti $\mathfrak{m}'$ terletak di atas $\mathfrak{m}$.
\begin{proof} 
Dibuktikan dengan cara serupa menggunakan teknik Artin. Saya agak lelah.
\end{proof} 
\end{comment}
\begin{example} 
Misalkan $R' = \mathbb{C}[x_1, \dots, x_n]/\mathfrak{p}$. Normalisasi Noether menyatakan
bahwa terdapat pemetaan injektif hingga $\mathbb{C}[y_1, \dots, y_a] \to R'$.
Klaimnya adalah
\[ \dim R'_{\mathfrak{m}} =a  \]
untuk setiap ideal maksimal $\mathfrak{m} \subset R'$. Kita telah siap membuktikan suatu
pernyataan yang sedikit lebih lemah. Khususnya (lihat di bawah untuk definisi
dimensi gelanggang nonlokal), berdasarkan proposisi kita
memperoleh klaim yang lebih lemah
\[ \dim R' = a,  \]
sebab dimensi gelanggang polinomial $\mathbb{C}[y_1, \dots, y_a]$ adalah $a$.
(\textbf{Saya rasa kita belum membuktikannya.})
\end{example} 



\section{Definisi dan karakterisasi lain dari dimensi}

\subsection{Karakterisasi topologis dimensi} Sekarang kita menginginkan suatu karakterisasi topologis
dimensi. Jadi, pertama-tama kita akan mempelajari bagaimana dimensi
berubah ketika kita melakukan operasi terhadap suatu modul. Misalkan $M$ suatu modul-$R$ terbangkitkan hingga atas gelanggang lokal
Noether $R$. Misalkan $x \in \mathfrak{m}$, dengan $\mathfrak{m}$ ideal
maksimalnya.
Kita dapat bertanya:
\begin{quote}
Apakah hubungan antara $\dim \supp M$ dan $\dim \supp M/xM$?
\end{quote}
Modul $M$ memetakan secara surjektif ke $M/xM$, sehingga kita mempunyai pertidaksamaan $\geq$. Namun kita memandang
dimensi sebagai banyaknya parameter yang diperlukan untuk mendeskripsikan
sesuatu. Banyaknya parameter seharusnya tidak berubah terlalu jauh ketika beralih dari
$M$ ke $M/xM$. Memang, sebagaimana dapat diperiksa,
\[ \supp M/xM = \supp M \cap V(x)  \]
dan mengiris $\supp M$ dengan ``hiperpermukaan'' $V(x)$ seharusnya menurunkan
dimensi sebesar satu. 


Karena itu, kita membuat dugaan berikut.
\begin{prediction}
\[ \dim \supp M/xM = \dim \supp M - 1.  \]
\end{prediction}
Jelas bahwa hal ini tidak selalu benar, misalnya jika $x$ bekerja sebagai nol pada $M$. Namun kita ingin
mengecualikan keadaan tersebut. 
Dalam kasus-kasus yang wajar, dugaan itu memang benar:

\begin{proposition} \label{dimdropsbyone}
Andaikan $x \in \mathfrak{m}$ bukan pembagi nol pada $M$. Maka 
\[ \dim \supp M/xM = \dim \supp M - 1.  \]
\end{proposition} 
\begin{proof} 
Untuk melihatnya, kita meninjau polinomial Hilbert. Pertimbangkan barisan eksak
\[ 0 \to xM \to M \to M/xM \to 0  \]
yang untuk setiap $t$ menghasilkan barisan eksak
\[ 0 \to xM/(xM \cap \mathfrak{m}^t M) \to M/\mathfrak{m}^t M \to M/(xM  +
\mathfrak{m}^t M) \to 0 . \]
Untuk $t$ besar, panjang unsur-unsur ini diberikan oleh polinomial Hilbert,
sebab suku di kanan adalah $M/xM \otimes_R R/\mathfrak{m}^t$. 
Kita mempunyai
\[ f_M^+(t) = f_{M/xM}^+(t) + \ell(xM/ (x M \cap \mathfrak{m}^t M), \quad t
\gg 0.  \]
Khususnya, $\ell( xM/ (xM \cap \mathfrak{m}^t M))$ merupakan polinomial dalam $t$.
Apa yang dapat kita katakan tentang polinomial itu? Perhatikan bahwa $xM \simeq M$ karena $x$ bukan pembagi nol.
Khususnya,
\[ xM / (xM \cap \mathfrak{m}^t M) \simeq M/N_t  \]
dengan
\[ N_t = \left\{a \in M: xa \in \mathfrak{m}^t M\right\} . \]
Khususnya, $N_t \supset \mathfrak{m}^{t-1} M$. Ini memberi tahu kita bahwa
$\ell(M/N_t) \leq \ell(M/\mathfrak{m}^{t-1} M) = f_M^+(t-1)$ untuk $t \gg 0$.
Dengan menggabungkannya dengan informasi di atas, kita memperoleh
\[ f_M^+(t) \leq f_{M/xM}^+(t) + f_M^+(t-1),   \]
yang mengimplikasikan bahwa $f_{M/xM}^+(t)$ sekurang-kurangnya sebesar selisih berturut-turut
$f_M^+(t) - f_M^+(t-1)$. Polinomial terakhir ini berderajat $\dim \supp M -1$. Khususnya,
$f_{M/xM}^+(t)$ berderajat paling sedikit $\dim \supp M -1 $. Ini memberikan
satu arah, bahkan arah yang sukar. Kita telah menunjukkan bahwa mengiris sesuatu dengan kodimensi satu
tidak menurunkan dimensinya terlalu jauh. 

Sekarang mari kita buktikan arah yang lain. Pada dasarnya, kita telah melakukannya sebelumnya melalui
Lema Artin--Rees. Kita mengetahui bahwa $N_t = \left\{a \in M: xa \in
\mathfrak{m}^t\right\}$. Lema Artin--Rees menyatakan bahwa terdapat konstanta
$c$ sedemikian sehingga $N_{t+c} \subset \mathfrak{m}^t M$ untuk setiap $t$. Karena itu,
$\ell(M/N_{t+c}) \geq \ell(M/\mathfrak{m}^t M) = f_M^+(t), t \gg 0$. Sekarang
ingat barisan eksak $0 \to M/N_t \to M/\mathfrak{m}^t M \to M/(xM +
\mathfrak{m}^t M) \to 0$. Dari sini kita melihat bahwa
\[ \ell(M/ \mathfrak{m}^t M) = \ell(M/N_t) + f_{M/xM}^+(t) \geq f_M^+(t-c) +
f_{M/xM}^+(t), \quad t \gg 0,  \]
yang mengimplikasikan bahwa
\[ f_{M/xM}^+(t) \leq f_M^+(t) - f_M^+(t-c),  \]
sehingga derajatnya harus turun. Kita memperoleh $\deg f_{M/xM}^+ < \deg f_{M}^+$.
\end{proof} 

Ini memberi kita suatu algoritma untuk menghitung dimensi modul-$R$ $M$. 
Pertama, masalahnya direduksi menjadi menghitung $\dim R/\mathfrak{p}$ untuk ideal prima $\mathfrak{p} \subset
R$. Kita dapat mengandaikan bahwa $R$ suatu daerah integral dan bahwa kita mencari
$\dim R$. Secara geometris, hal ini
bersesuaian dengan mengambil suatu komponen tak tereduksi dari $\spec R$.

Sekarang pilih sebarang $x
\in R$ sedemikian sehingga $x$ tak nol tetapi tidak invertibel. Jika unsur seperti itu tidak ada,
maka $R$ suatu medan dan berdimensi nol. Kemudian hitung $\dim R/x$
secara rekursif dan tambahkan satu.

Perhatikan bahwa algoritma ini tidak mengatakan apa pun mengenai polinomial Hilbert, dan hanya
membahas struktur ideal prima.

\subsection{Rekapitulasi}
Sebelumnya kita membahas teori dimensi. 
Ingat bahwa $R$ suatu gelanggang lokal Noether dengan ideal maksimal $\mathfrak{m}$,
dan $M$ suatu modul-$R$ terbangkitkan hingga. Kita dapat meninjau panjang $\ell(M/\mathfrak{m}^t M)$
untuk berbagai $t$; bagi $t \gg 0$, panjang ini merupakan fungsi polinomial. Derajat
polinomial ini disebut \textbf{dimensi} dari $\supp M$. 

\begin{remark} 
Jika $M = 0$, berdasarkan konvensi kita definisikan $\dim \supp M = -1$.
\end{remark} 

Sebelumnya, kita menunjukkan bahwa jika $M \neq 0$ dan $x \in \mathfrak{m}$ sedemikian sehingga $x$
bukan pembagi nol pada $M$ (yakni $M \stackrel{x}{\to} M$ injektif), maka 
\[ \boxed{ \dim \supp M/xM = \dim \supp M - 1.}\]
Dengan menggunakan hal ini, kita dapat memberikan suatu rekursi untuk menghitung dimensi. 
Untuk menghitung $\dim R = \dim \spec R$, perhatikan tiga sifat berikut:
\begin{enumerate}
\item $\dim R = \sup_{\mathfrak{p} \ \mathrm{ideal \ prima \ minimal}}
R/\mathfrak{p}$. Secara intuitif, ini mengatakan bahwa varietas yang merupakan gabungan
komponen-komponen tak tereduksi mempunyai dimensi yang sama dengan maksimum dimensi komponen-komponen tersebut.
\item $\dim R = 0$ jika $R$ suatu medan. Ini langsung dari definisi.
\item Jika $R$ suatu daerah integral dan $x \in \mathfrak{m} - \left\{0\right\}$, maka
$\dim R/(x) +1 = \dim R $. Ini langsung dari rumus berbingkai karena $x$ bukan pembagi nol.
\end{enumerate}

Ketiga sifat ini \emph{mencirikan secara tunggal} invarian dimensi. 

\textbf{Lebih tepatnya, jika
$d: \left\{\mathrm{gelanggang \ lokal \ Noether}\right\} \to \mathbb{Z}_{\geq 0}$
memenuhi ketiga sifat di atas, maka $d = \dim $. }
\begin{proof} 
Induksi pada $\dim R$. Jelas bahwa cukup membuktikannya ketika $R$ suatu daerah integral. 
Jika $R$ suatu medan, pernyataannya jelas; jika $\dim R>0$, syarat ketiga memungkinkan kita
mereduksi ke kasus yang dicakup hipotesis induksi (yakni turun satu tingkat).
\end{proof} 

Mari kita nyatakan ulang sifat 3 di atas:
\begin{quote}
3': Jika $R$ suatu daerah integral dan bukan medan, maka 
\[ \dim R = \sup_{x \in \mathfrak{m} - 0} \dim R/(x) + 1. \]
\end{quote}
Jelas bahwa 3' mengimplikasikan 3, dan dari argumen yang sama jelas bahwa 1, 2, dan 3'
mencirikan gagasan dimensi.

\subsection{Dimensi Krull} Sekarang kita akan mendefinisikan gagasan lain tentang
dimensi dan menunjukkan bahwa gagasan tersebut ekuivalen dengan yang lama dengan membuktikan bahwa ia
memenuhi aksioma-aksioma ini.

\begin{definition} 
Misalkan $R$ suatu gelanggang komutatif. Suatu \textbf{rantai ideal prima} dalam $R$ adalah barisan
hingga
\[ \mathfrak{p}_0 \subsetneq \mathfrak{p}_1 \subsetneq \dots \subsetneq
\mathfrak{p}_n.  \]
Rantai ini dikatakan mempunyai \textbf{panjang $n$.}
\end{definition} 

\begin{definition} 
\textbf{Dimensi Krull} dari $R$ adalah panjang maksimum rantai
ideal prima. Nilai ini mungkin saja $\infty$, tetapi segera akan kita lihat bahwa hal itu tidak dapat
terjadi jika $R$ lokal dan Noether.
\end{definition} 

\begin{remark} 
Untuk setiap rantai maksimal ideal prima $\left\{\mathfrak{p}_i, 0 \leq i \leq n\right\}$ (yakni yang tidak dapat diperpanjang),
$\mathfrak{p}_0$ harus merupakan ideal prima minimal dan $\mathfrak{p}_n$ suatu ideal maksimal.
\end{remark} 

\begin{theorem} 
Untuk gelanggang lokal Noether $R$, dimensi Krull dari $R$ ada dan sama
dengan $\dim R$ yang biasa.
\end{theorem}
\begin{proof} 
Kita akan menunjukkan bahwa dimensi Krull memenuhi aksioma-aksioma di atas. Untuk sementara,
tuliskan $\krdim$ untuk dimensi Krull.

\begin{enumerate}
\item Pertama, perhatikan bahwa $\krdim(R) = \max_{\mathfrak{p} \in R \
\mathrm{minimal}}  \krdim(R/\mathfrak{p})$. Hal ini karena setiap rantai ideal
prima dalam $R$ memuat suatu ideal prima minimal. Jadi setiap rantai ideal prima dalam $R$ dapat
dipandang sebagai suatu rantai dalam \emph{suatu} $R/\mathfrak{p}$, dan sebaliknya.
\item Kedua, kita perlu memeriksa bahwa $\krdim(R) = 0$ jika $R$ suatu medan. Hal ini
jelas, sebab hanya ada satu ideal prima.
\item Syarat ketiga menarik. Kita harus memeriksa bahwa untuk $(R,
\mathfrak{m})$ suatu gelanggang lokal yang merupakan
daerah integral, 
\[ \krdim(R) = \max_{x \in \mathfrak{m} - \left\{0\right\}} \krdim(R/(x)) + 1.  \]
Jika kita membuktikannya, berarti kita telah menunjukkan bahwa syarat 3' dipenuhi oleh
dimensi Krull. Dari argumen induktif di atas akan mengikuti bahwa $\krdim(R)
= \dim (R)$ untuk setiap $R$. 
Ada dua pertidaksamaan yang harus dibuktikan. Pertama, kita harus menunjukkan
\[ \krdim(R) \geq \krdim(R/x) +1, \quad \forall x \in \mathfrak{m} - 0.  \]
Andaikan $k = \krdim(R/x)$. Kita ingin menunjukkan bahwa terdapat rantai ideal
prima panjang $k+1$ dalam $R$. Misalkan $\mathfrak{p}_0 \subsetneq \dots
\subsetneq \mathfrak{p}_k$ suatu rantai panjang $k$ dalam $R/(x)$. Prapetanya
di $R$ memberikan rantai sejati ideal prima dalam $R$ yang panjangnya $k$, semuanya
memuat $(x)$ dan dengan demikian memuat $0$ secara sejati. Jadi, menambahkan nol menghasilkan rantai
ideal prima dalam $R$ yang panjangnya $k+1$. 

Sebaliknya, kita ingin menunjukkan bahwa jika terdapat rantai ideal prima dalam $R$ dengan
panjang $k+1$, maka terdapat rantai panjang $k$ dalam $R/(x)$ untuk suatu $x \in
\mathfrak{m} - \left\{0\right\}$. Tuliskan rantai panjang $k+1$ tersebut sebagai
\[ \mathfrak{q}_{-1} \subset \mathfrak{q}_0 \subsetneq \dots \subsetneq
\mathfrak{q}_k \subset R . \]
Jelas bahwa $\mathfrak{q}_0$ memuat suatu $x \in \mathfrak{m} - 0$. Maka rantai
$\mathfrak{q}_0 \subsetneq \dots \subsetneq \mathfrak{q}_k$ dapat
diidentifikasi dengan suatu rantai dalam $R/(x)$ untuk $x$ ini. Jadi untuk $x$ ini berlaku
$\krdim R \leq \sup \krdim R/(x) + 1$.
\end{enumerate}
\end{proof} 

Dengan demikian terdapat definisi kombinatorial bagi dimensi.

Secara geometris, misalkan $X = \spec R$ dengan $R$ suatu gelanggang afin atas $\mathbb{C}$ (suatu
gelanggang polinomial modulo suatu ideal). Maka $R$ mempunyai dimensi Krull $\geq k$ jika dan hanya jika terdapat
rantai subvarietas tak tereduksi dari $X$,
\[ X_0 \supset X_1 \supset \dots \supset X_k . \]
Pembenarannya akan dijumpai dalam \rref{subsectiondimension} di bawah.

\begin{remark}[\textbf{Peringatan!}] Misalkan $R$ suatu gelanggang lokal Noether berdimensi $k$. Ini
berarti terdapat rantai ideal prima panjang $k$, dan tidak ada rantai yang lebih
panjang. Jadi terdapat suatu rantai maksimal yang panjangnya $k$. Namun tidak semua
rantai maksimal dalam $\spec R$ mempunyai panjang $k$. 
\end{remark} 

\begin{example} 
Misalkan $R =( \mathbb{C}[X,Y,Z]/(XY,XZ))_{(X,Y,Z)}$. Sebagai
latihan, pembaca dapat memeriksa bahwa terdapat rantai-rantai maksimal yang
panjangnya bukan dua.

Terdapat \emph{daerah integral} lokal Noether yang lebih rumit dengan rantai-rantai maksimal
ideal prima yang tidak semuanya sama panjang. Contoh-contoh ini bukan contoh yang
dijumpai sehari-hari dan pasti rumit. Hal ini
tidak dapat terjadi untuk daerah integral terbangkitkan hingga atas suatu medan.
\end{example} 

\begin{example} 
Cara yang lebih sederhana agar tidak semua rantai maksimal mempunyai panjang yang sama adalah jika
$\spec R$ mempunyai dua komponen (dalam hal ini $R = R_0 \times R_1$ untuk gelanggang
$R_0, R_1$). 
\end{example} 


\subsection{Definisi lainnya lagi}
Mari kita mulai dengan memikirkan definisi modul. Ingat bahwa jika $(R,
\mathfrak{m})$ adalah
gelanggang lokal Noether, $M$ suatu modul-$R$ terbangkitkan hingga, dan $x \in \mathfrak{m}$
bukan pembagi nol pada $M$, maka
\[ \dim \supp M/xM = \dim \supp M -1.  \]

\begin{question} 
Bagaimana jika $x$ merupakan pembagi nol? 
\end{question} 

Pernyataan tersebut belum tentu benar (misalnya jika $x \in \ann(M)$). Meskipun demikian, kita mengklaim
bahwa bahkan dalam kasus ini:
\begin{proposition} 
Untuk setiap $x \in \mathfrak{m}$,
\[ \boxed{ \dim \supp M \geq \dim \supp M/xM \geq \dim \supp M -1 .}\]
\end{proposition} 
Batas atas bagi $\dim M/xM$ jelas karena $M/xM$ merupakan hasil bagi dari $M$. Batas
bawah lebih rumit. 

\begin{proof} 
Misalkan $N = \left\{a \in M: x^n a = 0 \ \mathrm{untuk \ suatu \ } n \right\}$. Kita dapat
membentuk barisan eksak
\[ 0 \to N \to M \to M/N \to 0.  \]
Misalkan $M'' = M/N$.
Berdasarkan konstruksi, $x$ bukan pembagi nol pada $M/N$. Kita klaim bahwa 
\[ 0 \to N/xN \to M/xM \to M''/xM'' \to 0  \]
juga eksak. Untuk itu kita hanya perlu melihat keeksakan pada bagian awal,
yakni keininjektifan $N/xN \to M/xM$. Jadi
kita perlu menunjukkan bahwa jika $a \in N$ dan $a \in xM$, maka $a \in x N$.

Untuk melihatnya, andaikan $a = xb$ dengan $b \in M$. Jika $\phi: M \to M''$, maka
$\phi(b) \in M''$ dianihilasi oleh $x$ sebab $x \phi(b) = \phi(bx) = \phi(a)$.
Ini berarti $\phi(b)=0$ karena $M'' \stackrel{x}{\to} M''$ injektif. Jadi
sebenarnya $b \in N$. Dengan demikian, $a \in xN$.

Dari keeksakan kita melihat bahwa (karena $x$ bukan pembagi nol pada $M''$)
\begin{align*} \dim M/xM & = \max (\dim M''/xM'', \dim N/xN) \geq \max(\dim M'' -1, \dim
N)\\ &  \geq \max( \dim M'', \dim N)-1  .  \end{align*}
Alasan untuk pernyataan terakhir adalah bahwa $\supp N/xN = \supp N$ sebab $N$ merupakan
modul torsi-$x$, dan dimensi hanya bergantung pada dukungan. Namun ruas kanan tepat sama dengan $\dim M -1$. 
\end{proof} 

Sebagai akibatnya, kita memperoleh:

\begin{proposition} 
$\dim \supp M$ adalah bilangan bulat minimal $n$ sedemikian sehingga terdapat unsur-unsur $x_1,
\dots, x_n \in \mathfrak{m}$ dan $M/(x_1 , \dots, x_n) M$ mempunyai panjang hingga. 
\end{proposition} 
Perhatikan bahwa $n$ selalu ada, sebab kita dapat mengambil sekumpulan pembangkit
ideal maksimal, dan $M/\mathfrak{m}M $ merupakan ruang vektor berdimensi hingga sehingga
mempunyai panjang hingga.
\begin{proof} 
Induksi pada $\dim \supp M$. Perhatikan bahwa $\dim \supp(M)=0$ jika dan hanya jika
polinomial Hilbert berderajat nol, yakni $M$ mempunyai panjang hingga atau $n=0$
(dengan $n$ didefinisikan seperti dalam pernyataan). 

Andaikan $\dim \supp M > 0$. \begin{enumerate}
\item Pertama-tama kita menunjukkan bahwa terdapat $x_1, \dots, x_{\dim M}$
sedemikian sehingga $M/(x_1, \dots, x_{\dim M})M$ mempunyai panjang hingga. 
Misalkan $M' \subset M$ submodul maksimal yang mempunyai panjang hingga. Terdapat
barisan eksak
\[ 0 \to M' \to M \to M'' \to 0  \]
dengan $M'' = M/M'$ tidak mempunyai submodul berpanjang hingga. Dalam hal ini, pada dasarnya kita dapat
mengabaikan $M'$ dan mengganti $M$ dengan $M''$. Alasannya adalah bahwa mengambil hasil bagi
oleh $M'$ tidak memengaruhi $n$ maupun dimensi. 

Jadi, ganti $M$ dengan
$M''$, sehingga kita dapat mengandaikan bahwa $M$ tidak mempunyai submodul berpanjang hingga. Khususnya,
$M$ tidak memuat salinan $R/\mathfrak{m}$, yakni $\mathfrak{m}
\notin \ass(M)$. 
Menurut penghindaran ideal prima, ini berarti terdapat $x_1 \in \mathfrak{m}$ yang bekerja sebagai
bukan pembagi nol pada $M$. Jadi,
\[ \dim M/x_1M = \dim M -1.  \]
Hipotesis induksi menyatakan bahwa terdapat $x_2, \dots, x_{\dim M}$ sedemikian sehingga
$$(M/x_1 M)/(x_2, \dots, x_{\dim M}) (M/xM) \simeq M/(x_1, \dots, x_{\dim M})M $$
mempunyai panjang hingga. Ini membuktikan klaim tersebut.
\item Sebaliknya, andaikan $M/(x_1, \dots, x_n)M$ mempunyai panjang hingga.
Kita klaim bahwa $n \geq \dim M$. Hal ini mengikuti dari hasil sebelumnya
bahwa mengambil hasil bagi oleh satu unsur dapat menurunkan dimensi paling banyak
$1$. Dengan menerapkannya secara rekursif, dan menggunakan fakta bahwa $\dim$ dari modul
berpanjang hingga adalah nol, diperoleh
\[ 0 = \dim M/(x_1 , \dots, x_n )M \geq \dim M -n. \]
\end{enumerate}
\end{proof} 


\begin{corollary} 
Misalkan $(R, \mathfrak{m})$ suatu gelanggang lokal Noether. Maka $\dim R$ sama dengan $n$ minimal
sedemikian sehingga terdapat $x_1, \dots, x_n \in R$ dengan $R/(x_1, \dots, x_n) R$
bersifat Artin. Atau, secara ekuivalen, sedemikian sehingga $(x_1, \dots, x_n)$ memuat suatu pangkat
$\mathfrak{m}$.
\end{corollary}


\begin{remark} 
Di sini terlihat jelas bahwa dimensi $R$ paling besar sama dengan dimensi
pembenaman. Di sini kita tidak menuntut pembangkitan ideal maksimal, melainkan
hanya suatu ideal yang memuat pangkat ideal maksimal tersebut.
\end{remark} 
\lecture{11/5}

Kita telah membahas dimensi. Misalkan $R$ suatu gelanggang lokal Noether dengan
ideal maksimal $\mathfrak{m}$. Seperti yang dikatakan dalam kuliah sebelumnya, $\dim R$ dapat dicirikan oleh:
\begin{enumerate}
\item Bilangan minimal $n$ sedemikian sehingga terdapat suatu ideal $n$-primer yang dibangkitkan oleh $n$
unsur $x_1, \dots, x_n \in \mathfrak{m}$. Artinya, titik tertutup
$\mathfrak{m}$ dari
$\spec R$ dipotong secara \emph{teoretis-himpunan} oleh irisan $\bigcap
V(x_i)$. Ini merupakan salah satu cara mengatakan bahwa titik tertutup dapat didefinisikan oleh $n$
parameter. 
\item Bilangan \emph{maksimal} $n$ sedemikian sehingga terdapat rantai ideal prima
\[ \mathfrak{p}_0 \subset \mathfrak{p}_1 \subset \dots \subset \mathfrak{p}_n. \]
\item Derajat polinomial Hilbert $f^+(t)$, yang sama dengan
$\ell(R/\mathfrak{m}^t)$ untuk $t \gg 0$.
\end{enumerate}


\subsection{Hauptidealsatz Krull}


Misalkan $R$ suatu gelanggang lokal Noether.
Pernyataan berikut sekarang jelas dari apa yang telah kita tunjukkan:

\begin{theorem} \label{hauptv1}
$R$ berdimensi $1$ jika dan hanya jika terdapat bukan pembagi nol $x \in \mathfrak{m}$ sedemikian sehingga
$R/(x)$ bersifat Artin.
\end{theorem} 



\begin{remark} 
Misalkan $R$ suatu daerah integral. Kita mengatakan bahwa ideal prima tak nol $\mathfrak{p} \subset R$
\textbf{bertinggi satu} jika $\mathfrak{p}$ minimal di antara ideal-ideal prima
yang memuat suatu $x \in R$ tak nol. 

Menurut Hauptidealsatz Krull, $\mathfrak{p}$ bertinggi satu \textbf{jika
dan hanya jika $\dim R_{\mathfrak{p}} = 1$.}
\end{remark} 


Kita dapat memperumum gagasan mengenai $\mathfrak{p}$ sebagai berikut.
\begin{definition} 
Misalkan $R$ suatu gelanggang Noether (tidak harus lokal), dan $\mathfrak{p} \in
\spec R$. Kita definisikan \textbf{tinggi} dari $\mathfrak{p}$, yang dilambangkan
$\het(\mathfrak{p})$, sebagai $\dim R_{\mathfrak{p}}$.
Kita mengetahui bahwa nilai ini adalah panjang suatu rantai maksimal ideal prima dalam
$R_{\mathfrak{p}}$. Jadi ini merupakan panjang maksimal rantai ideal prima dari $R$, 
\[ \mathfrak{p}_0 \subset \dots \subset \mathfrak{p}_n = \mathfrak{p}  \]
yang berakhir pada $\mathfrak{p}$. Inilah asal istilah ``tinggi''.
\end{definition} 

\begin{remark} 
Kadang-kadang tinggi disebut \textbf{kodimensi}. Ini bersesuaian dengan
kodimensi dalam $\spec R$ dari subhimpunan tertutup tak tereduksi yang bersesuaian dalam
$\spec R$.
\end{remark} 

\begin{theorem}[Hauptidealsatz Krull] Misalkan $R$ suatu gelanggang Noether dan $x
\in R$ bukan pembagi nol. Jika $\mathfrak{p} \in \spec R$ minimal di atas $x$,
maka $\mathfrak{p}$ bertinggi satu.
\end{theorem} 
\begin{proof} 
Langsung dari \cref{hauptv1}.
\end{proof} 

\begin{theorem}[Artin-Tate]
Misalkan $A$ suatu daerah integral Noether. Maka pernyataan-pernyataan berikut ekuivalen:
\begin{enumerate}
\item Terdapat $f \in A-\left\{0\right\} $ sedemikian sehingga $A_f$ suatu medan.
\item $A$ mempunyai berhingga banyak ideal maksimal dan berdimensi paling besar 1.
\end{enumerate}
\end{theorem} 
\begin{proof} Kita mengikuti \cite{EGA}.

Pertama, andaikan terdapat $f$ sedemikian sehingga $A_f$ suatu medan. 
Maka semua ideal prima tak nol dari $A$ memuat $f$. 
Kita perlu menyimpulkan bahwa $A$ berdimensi $\leq 1$. Tanpa mengurangi keumuman,
kita dapat mengandaikan bahwa $A$ bukan medan.

Terdapat berhingga banyak ideal prima $\mathfrak{p}_1,\dots, \mathfrak{p}_k$ yang
minimal di atas $f$; semuanya bertinggi satu. Klaimnya adalah bahwa setiap ideal maksimal $A$ mempunyai
bentuk ini. Andaikan $\mathfrak{m}$ maksimal dan bukan salah satu dari $\mathfrak{p}_i$.
Menurut penghindaran ideal prima, terdapat $g \in \mathfrak{m}$ yang
tidak berada dalam satu pun $\mathfrak{p}_i$. Suatu ideal prima minimal $\mathfrak{P}$ di atas $g$ bertinggi
satu, sehingga berdasarkan asumsi kita memuat $f$. Namun dengan demikian ideal itu merupakan salah satu
$\mathfrak{p}_i$; ini suatu kontradiksi karena $g \in \mathfrak{P}$.
\end{proof} 


\subsection{Catatan lebih lanjut}

Kita dapat menyatakan kembali gagasan-gagasan sebelumnya dalam istilah dimensi.
\begin{remark} 
Suatu gelanggang Noether berdimensi nol jika dan hanya jika $R$ bersifat Artin. Memang,
$R$ berdimensi nol jika dan hanya jika semua ideal prima bersifat maksimal.
\end{remark} 


\begin{remark} 
Suatu daerah integral Noether berdimensi nol jika dan hanya jika merupakan medan. Memang, dalam hal ini
$(0)$ maksimal.
\end{remark} 

\begin{remark} 
$R$ berdimensi $\leq 1$ jika dan hanya jika setiap ideal prima nonminimal dari $R$
bersifat maksimal. Artinya, tidak ada rantai yang panjangnya $\geq 2$.
\end{remark} 

\begin{remark} 
Suatu daerah integral (Noether) $R$ berdimensi $\leq 1$ jika dan hanya jika setiap ideal prima tak nol
bersifat maksimal.
\end{remark} 

Khususnya,
\begin{proposition} 
$R$ merupakan gelanggang Dedekind jika dan hanya jika merupakan daerah integral Noether yang tertutup secara integral dan berdimensi
$1$. 
\end{proposition} 
	
	
